% !PRINTED 29
\weberpage{29}

Die Coëfficienten der Functionen $Q$ und $C$ sind aus den
Coëfficienten $a$ und $b$ durch Addition, Subtraction, Multiplication
und Theilung zusammengesetzt. Im Nenner kommen aber nur Potenzen von
$b_0$ vor, und wenn also $b_0=1$ ist, so sind die Coëfficienten von
$Q$ und $C$ ganze Functionen der $a$ und $b$. Die höchste Potenz von
$b_0$, die im Nenner auftreten kann, ist die
$(m-n+1)$\textsuperscript{te}, da in der Kette (4) in jeder folgenden
Gleichung im Nenner einmal der Factor $b_0$ hinzukommt. Es kann aber
in besonderen Fällen die höchste Potenz von $b_0$ in allen Nennern
eine niedrigere sein.

Nehmen wir z. B. für $A$ eine Function dritten Grades
\begin{equation}
 f(x)=a_0x^3+a_1x^2+a_2x+a_3
\tag{7}
\end{equation}
und für $B$ die sogenannte erste Derivirte von $f(x)$, die vom zweiten
Grade ist,
\begin{equation}
 f'(x)=3a_0x^2+2a_1x+a_2,
\tag{8}
\end{equation}
so erhält man:
\begin{equation}
 Q={1\over 3}x+{a_1\over 9a_0},
\tag{9}
\end{equation}
\begin{equation}
 C={6a_0a_2-2a_1^2\over 9a_0}x
   +{9a_0a_3-a_1a_2\over 9a_0}.
\tag{10}
\end{equation}

Wie man die in den Gleichungen (4) vorgeschriebene Rechnung
zweckmässig anordnet, darf hier aus den Elementen als bekannt
vorausgesetzt werden. Wir machen auf die Analogie aufmerksam, die
zwischen dieser Rechnung und der Division im dekadischen Zahlsystem
besteht. Eine Function $f(x)$ stellt eine dekadisch geschriebene Zahl
dar, wenn die Coëfficienten $a_0,a_1,\ldots$ ganze Zahlen zwischen Null
(einschliesslich) und $10$ (ausschliesslich) sind, und $x=10$ gesetzt
wird. Lässt man in den Coëfficienten auch Zahlen zu, die grösser als
$10$ sind, so kann man eine Zahl auf verschiedene Arten durch $f(x)$
darstellen. Man wendet dies bei dem Divisionsverfahren an, um auf die
einfachste Weise in den Resultaten gebrochene und negative
Coëfficienten zu vermeiden.
